%----------------------------------------------------------------------
\section{The Attribute Grammar}\label{sec:grammar}\footnote{Section-level Toulmin. \emph{Move}: fix the RISC-style atom inventory and the base productions that combine atoms into plans. \emph{[DAF] comparison}: [DAF] builds sets from discriminators via $\kappa$-evolution and residues; we build plans from atoms via productions and cost attributes. \emph{Residue instance}: each production's cost vector is the residue of the difference between its children's sum-of-costs and the combined cost after fusion, time-sharing, or encoding change.}
%----------------------------------------------------------------------

\subsection{The atom inventory}\label{sec:atoms}

The grammar's terminals are the \emph{atoms}: primitive
kernels corresponding to one hardware dispatch each. The
inventory is intentionally small. Every compound behavior
is built from these terminals via the productions of
Section \ref{sec:productions}.

\begin{definition}[Atom]\label{def:atom}
  An \emph{atom} is a quadruple $(\mathrm{id},
  \mathrm{sig}, \pi_0, \mathrm{body})$ where
  $\mathrm{id}$ is a name in the set
  \begin{center}
  \begin{tabular}{ll}
    $\oplus$ & pointwise XOR in encoding $E$ \\
    $\wedge_{\mathrm{tile}}$ & GF(2) AND in encoding $E$ \\
    $\mathrm{pcnt}_2$ & popcount mod 2 (GF(2) reduce) \\
    $\zeta_{\mathrm{pass}}$ & one butterfly pass of the zeta transform \\
    $\partial_{\mathrm{gather}}$ & boundary gather--XOR at fixed $n$ \\
    $\wedge_{\mathrm{yates}}$ & ranked-zeta wedge step \\
    $\mathrm{rot}_{2 \times 2}$ & perspective rotation tile \\
    $\mathrm{cd}_{2 \times 2}$ & Cayley--Dickson tile \\
    $\mathrm{syn}_{7 \times 3}$ & Hamming syndrome tile \\
    $\mathrm{fano}_3$ & Fano-line predicate \\
    $\mathrm{enc}_{\Lambda \to \mathrm{RM}}$ & encode $\Lambda \to \mathrm{RM}$ ($n$ butterflies) \\
    $\mathrm{dec}_{\mathrm{RM} \to \Lambda}$ & decode $\mathrm{RM} \to \Lambda$ (identical to encode over $\GF(2)$) \\
    $\mathrm{enc}_{\Lambda \to \mathrm{Ham}}$ & encode $\Lambda \to $ Hamming compressed \\
    $\mathrm{dec}_{\mathrm{Ham} \to \Lambda}$ & decode Hamming $\to \Lambda$ \\
    $\mathrm{gather}$, $\mathrm{scatter}$ & stride-regular memory reshape \\
  \end{tabular}
  \end{center}
  $\mathrm{sig}$ is the atom's signature (input and
  output encoding, grade budget, tile size, batch
  class); $\pi_0$ is its baseline cost vector drawn from
  the hardware profile; $\mathrm{body}$ is the Pallas
  kernel text (Appendix \ref{app:kernels}).
\end{definition}

\begin{remark}
  The atom inventory is closed under [DAF]'s algebra in
  the sense that every CSTZ operation from \S\S3--9 of
  [DAF] decomposes into finitely many atoms from the
  list above. Section \ref{sec:discussion} gives the
  explicit mapping. No atom is marked as non-linear
  except $\mathrm{pcnt}_2$, which reduces an
  $n$-bit input to a single bit.
\end{remark}

\subsection{The attribute domain}\label{sec:attributes}

Each nonterminal (= partial plan) carries an attribute
tuple
\begin{equation}
  \bar{a} = (E, r, d, k; \bar{c}) \in
  \mathcal{E} \times \mathbb{N} \times \mathbb{N}
  \times \mathbb{N} \times \mathbb{N}^4,
\end{equation}
where $E$ is the output encoding (an element of the
encoding lattice $\mathcal{E}$ of
Section \ref{sec:encoding}), $r$ the grade budget, $d$
the natural tile size of the output, $k$ the batch
class, and $\bar{c}$ the cost vector.

Two plans with identical $(E, r, d, k)$ are
comparable: their $\bar{c}$ values partially order
them. The Pareto-nondominated subset of a set of
plans with identical semantic attributes is what the
parser carries forward; Section \ref{sec:portfolios}
shows how to extract it.

\subsection{Base productions}\label{sec:productions}

\begin{definition}[Plan]\label{def:plan}
  A \emph{plan} is a derivation tree rooted at the
  workload's output nonterminal, whose leaves are
  atoms and whose internal nodes are productions. The
  grammar consists of the productions below.
\end{definition}

Three classes of productions operate on plans. The
first (base) corresponds directly to atoms; the second
class (linear) fuses chains of $\GF(2)$-linear atoms;
the third (compound) introduces spatial fusion
(Section \ref{sec:fusion}) or temporal time-sharing
(Section \ref{sec:time-sharing}).

\subsubsection*{Base productions}

\begin{center}
\small
\begin{tabular}{l}
\verb|Plan(E, r, d, k) ::= Atom[E, r, d, k]|\\[2pt]
\verb|Plan(E2, r, d, k) ::= enc[E1 -> E2] Plan(E1, r, d, k)|\\[2pt]
\verb|Plan(E, r-1, d, k) ::= partial[E] Plan(E, r, d, k)|\\[2pt]
\verb|Plan(E, r1+r2, d, k) ::= wedge[E] Plan(E, r1, d, k) Plan(E, r2, d, k)|
\end{tabular}
\end{center}

Each base production's cost is the concatenation of
its children's costs plus the atom's $\pi_0$, under
the \emph{serial} cost-combination rule:
\begin{align}
  \VRAM &= \max_i \VRAM_i + \VRAM_{\mathrm{atom}} \\
  \SMEM &= \max_i \SMEM_i + \SMEM_{\mathrm{atom}} \\
  \REG  &= \max_i \REG_i  + \REG_{\mathrm{atom}}  \\
  \LAT  &= \sum_i \LAT_i + \LAT_{\mathrm{atom}}
\end{align}
which captures that base productions run their
children in sequence: latencies add, but peak resources
are not simultaneously live (one child's resources are
released before the next child begins).

\subsubsection*{Linearity-fusion productions}

\begin{center}
\small
\begin{tabular}{l}
\verb|Plan(E, r, d, k+k') ::= xor_fused Plan(E, r, d, k) Plan(E, r, d, k')|
\end{tabular}
\end{center}

where both children end with an $\oplus$ atom on the
same $(E, r, d)$. This exploits
Proposition \ref{prop:encoding-commute}: chained XORs
combine into one by associativity, amortizing the
kernel-launch cost across children.

\subsection{Soundness}

\begin{theorem}[Grammar soundness]\label{thm:soundness}
  \footnote{Class \textsf{A}. Warrant: structural
  induction on the grammar productions. Qualifier:
  assumes each atom's Pallas body is faithful to its
  stated $\GF(2)$ semantics; this is the subject of
  property-based equivalence testing rather than of
  grammar-level proof. Finite-$\kappa$ valid.}
  Let $P$ be any plan derivable under the grammar
  from a workload $W$. Executing $P$ on a stated
  hardware profile produces the same $\GF(2)$ output
  as executing $W$ in the reference interpreter of
  [DAF].
\end{theorem}

\begin{proof}
  Induction on the derivation of $P$. Base case: a
  single atom is semantically identical to its
  reference interpretation by the atom's
  specification. Inductive case: each production's
  semantic rule is a $\GF(2)$-linear identity
  (XOR fusion, encoding commutation
  (Proposition \ref{prop:encoding-commute}), grade
  projection), compound productions
  (Sections \ref{sec:fusion}--\ref{sec:time-sharing})
  have their own correctness propositions
  (Propositions \ref{prop:fused-correctness},
  \ref{prop:xor-bridge}), and the compositional
  structure of the grammar composes these locally
  sound rules into a globally sound derivation.
\end{proof}

\begin{remark}[Reference interpretation]
  The reference interpretation of $W$ is the naive
  sequential execution of each atom in its specified
  encoding, with explicit materialization at every
  atom boundary. This is the ``oracle'' used in the
  property-based testing of item S4 in the reading
  guide: random workloads generated at small $n$ are
  executed twice, once by the reference interpreter
  and once by a planned schedule, and their $\GF(2)$
  outputs are XOR-compared for equality.
\end{remark}

\subsection{Well-formedness}

\begin{definition}[Well-formed grammar]\label{def:wf-grammar}
  The grammar is \emph{well-formed} if
  \begin{enumerate}[nosep,label=(\roman*)]
    \item every production preserves the semantic
      attributes $(E, r, d)$ exactly (i.e.\ the output
      encoding, grade budget, and tile size are
      functions of the children's attributes only, not
      of cost);
    \item no production is left-recursive on a
      non-linear atom (ensures the Earley parser of
      Section \ref{sec:parsing} terminates);
    \item the admissible heuristic $h^*$ of
      Section \ref{sec:parsing} is monotone
      (Section \ref{sec:parsing}, Proposition
      \ref{prop:heuristic-monotone}).
  \end{enumerate}
\end{definition}

\begin{proposition}[Base grammar is well-formed]
  \label{prop:base-wf}
  The productions defined above satisfy conditions
  (i)--(iii) of Definition \ref{def:wf-grammar}.
\end{proposition}

\begin{proof}
  (i) is direct from the production signatures above.
  (ii) holds because the only non-linear atom is
  $\mathrm{pcnt}_2$, which has arity 1 and therefore
  cannot be left-recursive. (iii) is deferred to
  Proposition \ref{prop:heuristic-monotone} in
  Section \ref{sec:parsing}.
\end{proof}

\subsection{Commutativity within productions}

\begin{proposition}[XOR-fusion commutativity]
  \label{prop:xor-commute}
  \footnote{Class \textsf{A}. Warrant: XOR is
  commutative and associative. Finite-$\kappa$ valid.}
  Under the \texttt{xor\_fused} production, the
  ordering of children is irrelevant: any permutation
  produces an equivalent plan.
\end{proposition}

\begin{proof}
  By commutativity and associativity of $\oplus$ in
  $\GF(2)$. The cost vector is symmetric under
  permutation because each component's combination
  rule (max or sum) is commutative.
\end{proof}

This innocuous-looking proposition licenses the
parser to avoid enumerating permutations of fused
XOR subplans, shrinking the chart by a factor of
$k!$ for a $k$-way fusion.

\subsection{Relationship to [DAF]}

The grammar structure mirrors [DAF]'s discrimination
regime:

\begin{center}
\small
\begin{tabular}{l|l}
\textbf{[DAF] construct} & \textbf{This paper's analogue} \\
\hline
Discriminator $d \in \bigwedge^1 V$ & Atom $a \in \mathcal{A}$ \\
$\kappa$-evolution & Production rule firing in the chart \\
Residue $R(S)$ & Cost vector difference \\
Membership profile $(\tau, \sigma)$ & Semantic attribute $(E, r)$ \\
Operationalist axiom & Execution-operationalist (Def.\ \ref{def:exec-operationalist}) \\
Sheafification & Chart closure under completer \\
SPPF & Plan portfolio (\S\ref{sec:portfolios}) \\
\end{tabular}
\end{center}

Each row is a functorial correspondence rather than
a strict identity, but the analogy is close enough
that this paper's machinery can be described
informally as ``applying [DAF]'s sheaf-theoretic
pattern to the compiler's scheduling problem.''
